\documentclass[12pt,a4paper]{article}

\usepackage{fontspec}
\usepackage{amsmath, amssymb, amstext}
\usepackage{graphicx}
\usepackage{geometry}
\usepackage{xcolor}
\usepackage{fancyhdr}

\geometry{margin=1in}

\pagestyle{fancy}
\setlength{\headheight}{15pt}

\lhead{کوییز درس فیزیک مدرن}
\chead{دانشکده فیزیک دانشگاه صنعتی خواجه نصیرالدین طوسی}
\rhead{۱۴۰۵/۰۴/۰۲}

\usepackage{xepersian}
\settextfont{B Nazanin}

\begin{document}
	
	\begin{tabular}{|p{7cm}|p{7cm}|}
		\hline
		\textbf{نام و نام خانوادگی:} & \textbf{شماره دانشجویی:} \\
		\hline
	\end{tabular}
	
	\vspace{10pt}
	
	\section*{سوالات}
	
	\begin{enumerate}
		
		\item 
		
	
	
	طول موج دوبروي ذرات زير را محاسبه کنيد:
		
	\begin{enumerate}
		\item[(الف)] یک پروتون با انرژی \(5\,\mathrm{MeV}\)
		\item[(ب)] یک الکترون با انرژی \(60\,\mathrm{GeV}\)
		\item[(ج)] یک الکترون با سرعت \(2.10 \times 10^6\,\mathrm{m/s}\)
	\end{enumerate}
		
		\item
		
		
		با در نظر گرفتن یک توزیع ساده از عدد موج ها در بازه ای به عرض \(\Delta k\) حول مقدار \(k_0\)، فرض کنید:
		
		\[
		A(k) =
		\begin{cases}
			A_0 & k_0 - \dfrac{\Delta k}{2} \le k \le k_0 + \dfrac{\Delta k}{2} \\
			0 & \text{\lr{otherwise}}
		\end{cases}
		\]
		
		و تابع موج به صورت زیر تعریف شود:
		
		\[
		y(x) = \int A(k)\cos(kx)\,dk
		\]
		
		(الف) با انجام انتگرال نشان دهید:
		
		\[
		y(x) = \frac{2A_0}{x}\,\sin\left(\frac{\Delta k\, x}{2}\right)\cos(k_0 x)
		\]
		
		(ب) نشان دهید این تابع یک بسته موج \lr{(wave packet)} است که شامل یک نوسان سریع با فرکانس \(k_0\) و یک پوش آرام با پهنای وابسته به \(\Delta k\) می باشد.
		
	\end{enumerate}
	
	\vspace{10pt}
	\textbf{موفق باشید.}
	
\end{document}